\section{Additional results}
\label{app:more}

\subsection{The deflated Sharpe ratio's hurdle}
\label{app:dsr}
Formally, the rule tests $H_0\colon SR\le SR_0$, where $SR_0$ is the expected maximum of $N$ skill-less trials, a harder hypothesis than ``the best strategy has no skill''. If $1-\mathrm{DSR}$ were a calibrated $p$-value for no skill, a value of at least 0.95 would occur in about 5\% of skill-less searches; it occurred in at most \nz{v1.size.deflated_sharpe.correlated_trials}\%, so at its 0.95 point the rule is strongly conservative for that hypothesis. The difference in hurdles is large:
\begin{itemize}
  \item Under independence and normality, the one-sided 5\% critical value for the best of $N=\nz{trials}$ trials is about \nz{dsr.sidak_crit} standard errors.
  \item The DSR's benchmark is about \nz{dsr.exp_max} cross-trial standard deviations (the approximation of \citet{bailey2014pseudo}; the exact expected maximum of 20 standard normals is \nz{dsr.exp_max_exact}), and the rule asks for a further \nz{dsr.single_crit} standard errors on top of that.
  \item The cross-trial standard deviation is computed from all \nz{trials} strategies, including the best. In the (unregistered) v0-design robustness run it averaged \nz{dsr.sd_ratio_null} times the i.i.d.\ standard error under the null, a hurdle of about \nz{dsr.hurdle_null} standard errors; under the alternative the skilled strategy widens the dispersion to \nz{dsr.sd_ratio_skill} times the standard error and the hurdle to about \nz{dsr.hurdle_skill}: a skilled strategy raises its own benchmark.
  \item When strategies are strongly correlated the dispersion shrinks; below \nz{dsr.sd_ratio_crossover} times the standard error the DSR's hurdle falls under the \v{S}id\'ak value. In the correlated family it was \nz{dsr.sd_ratio_null_corr}, a hurdle of about \nz{dsr.hurdle_null_corr}, which is why the DSR rejects most there.
\end{itemize}
The ranking the DSR induces is better than its raw power suggests: size-adjusted, its power under i.i.d.\ returns is \nz{v1.adj2.deflated_sharpe.iid_normal}\% against \nz{v1.adj2.luck_trials.iid_normal}\% for the \v{S}id\'ak test, and across families it trails by \nz{v1.dsr_adj_gap_range} points. Under i.i.d.\ returns the DSR's raw power falls \nz{v1.dsr_gap.total} points short of the \v{S}id\'ak test's size-adjusted power; \nz{v1.dsr_gap.threshold} of those points come from the 0.95 threshold (the DSR's own size-adjusted power is that much higher than its raw power), and the remaining \nz{v1.dsr_gap.statistic} from the statistic, whose benchmark moves with the dispersion of the whole search.

\subsection{Resolution matters for the bootstrap}
\label{sec:bootres}
A bootstrap $p$-value with $B$ resamples is a Monte Carlo estimate, and a test built on it loses power as $B$ falls \citep{davidson2000bootstrap}. With a \v{S}id\'ak adjustment the loss is sharp: with 400 resamples the smallest attainable adjusted $p$-value is \nz{boot.min_p.400}, just below the level, so the test rejects only when no resample reaches the observed Sharpe ratio. Scoring the same \nz{br.reps} searches per cell with 400 and with 1,999 resamples in an unregistered paired run (Table~\ref{tab:bootres}), the bootstrap's power under i.i.d.\ normal returns rises from \nz{br.power.b400.iid_normal}\% to \nz{br.power.b1999.iid_normal}\% (the \v{S}id\'ak $t$ test: \nz{br.power.let.iid_normal}\%), while its size moves from \nz{br.size.b400.iid_normal}\% to \nz{br.size.b1999.iid_normal}\%. Comparisons of resampling tests with analytical ones should state the number of resamples.

\subsection{The Harvey--Liu haircut read one-sided is the \v{S}id\'ak test}
For a positive Sharpe ratio, the two-sided $p\le 0.05$ of \citet{harvey2015backtesting} is a one-sided test at 2.5\%, which is why the haircut with its authors' convention sits near half the level (\nz{v1.size.haircut_bonferroni.iid_normal}\% under i.i.d.\ normal returns). Read one-sided, the haircut and the luck-equivalent-trials test use the same Student $t$ tail $p_1$ of the same statistic; Bonferroni rejects when $Np_1\le0.05$ and \v{S}id\'ak when $1-(1-p_1)^N\le0.05$, and the two disagree only when $p_1$ lies between \nz{bonf.p1} and \nz{sidak.p1}. Their sizes differ by at most \nz{v1.let_hc_maxdiff} points in any family. This is a known equivalence, not a finding; its practical consequence is that comparisons of these corrections at unmatched sides measure the choice of sides.

\subsection{The bootstrap's own variance}
\label{sec:inherit}
The bootstrap's critical value carries the same estimate. Under the stationary bootstrap the conditional variance of a resampled mean is the Politis--Romano estimate, so every block-bootstrap test, studentized or not, inherits that noisy, downward-biased estimate through its critical value, and the maximum over many strategies averages part of the noise out. Experiment E of v2b, registered after the v2 results, measures how much this matters under i.i.d.\ normal returns (Table~\ref{tab:k-block}). With one strategy, the Reality Check rejected \nz{v2b.E.size.rc.k1.b1}\% of skill-less searches at block length 1, \nz{v2b.E.size.rc.k1.b8}\% at 8 and \nz{v2b.E.size.rc.k1.b22}\% at 22 (P11 \nz{v2b.P11.held}). At block length 8 the effect is small: the Reality Check and the bootstrap-$t$ stayed within the Monte Carlo band for every number of strategies from 1 to 100 (\nz{v2b.E.range.rc.b8}\% and \nz{v2b.E.range.spa_c_boot_t_ge.b8}\%), and P12, which expected the Reality Check's size with one strategy to exceed its size with a hundred by at least half a point at block length 8, failed: the difference was \nz{v2b.P12.diff} points. At block length 22 both were liberal with up to 20 strategies (\nz{v2b.E.range_k20minus.rc.b22}\% and \nz{v2b.E.range_k20minus.spa_c_boot_t_ge.b22}\%), and the maximum over 50 or 100 strategies damped the effect (\nz{v2b.E.size.rc.k50.b22}\% and \nz{v2b.E.size.rc.k100.b22}\% for the Reality Check). The fixed-variance test compounds the two effects. At block length 8 its size rose from \nz{v2b.E.size.spa_c_ge.k1.b8}\% with one strategy to \nz{v2b.E.size.spa_c_ge.k100.b8}\% with a hundred, and at block length 22 from \nz{v2b.E.size.spa_c_ge.k1.b22}\% to \nz{v2b.E.size.spa_c_ge.k100.b22}\%; at block length 1, where its variance is the sample variance, it stayed at \nz{v2b.E.range.spa_c_ge.b1}\%. With twenty strategies and block length 8 it rejected \nz{v2b.E.size.spa_c_ge.k20.b8}\% of searches, a fourth replication of the main result on fresh seeds (P10 \nz{v2b.P10.held}).

Two v2 observations that pointed to the number of strategies were not replicated. With five i.i.d.\ strategies at block length 8, experiment B had found \nz{v2.B.size.rc.iid_normal.n504_k5}\% for the Reality Check and \nz{v2.B.size.spa_c_boot_t.iid_normal.n504_k5}\% for the bootstrap-$t$; experiment E, on fresh seeds, found \nz{v2b.E.size.rc.k5.b8}\% and \nz{v2b.E.size.spa_c_boot_t_ge.k5.b8}\%. In the correlated and cluster families, however, the Reality Check and the bootstrap-$t$ rejected \nz{corrclus.range}\% of skill-less searches across the runs, at or just above the upper edge of the band, which the number of independent strategies alone does not explain. P2's failure in v2, the sample-standard-deviation version at \nz{v2.A.size.spa_c_sd.block_cluster}\% in the cluster family, belongs to that pattern.

\subsection{Re-studentizing with a block variance}
\label{sec:block-student}
Re-studentizing every resample by its own standard deviation ignores serial dependence, and under AR(1) returns that bootstrap-$t$ was liberal (\nz{v2.A.size.spa_c_boot_t.ar1}\% in v2, \nz{v2b.F.size.spa_c_boot_t_ge.ar1}\% in v2b). Experiment F re-studentized each resample instead by its natural block variance---the squared sums, within each resampled block, of the deviations from the resample mean, averaged over the $T$ periods \citep{gotze1996second}---while the observed statistic kept the Politis--Romano studentization (Table~\ref{tab:block-student}). Under AR(1) returns this brought the size to \nz{v2b.F.size.spa_c_nb_ge.ar1}\% (P14 \nz{v2b.P14.held}), and outside the AR(1) family it never exceeded 6.5\% (P13 \nz{v2b.P13.held}). But it over-corrects: outside the AR(1) family its size was \nz{v2b.F.size_range_nonar.spa_c_nb_ge}\%, below the band in \nz{v2b.F.below_band_nonar_text.spa_c_nb_ge} families, and experiment E shows it growing more conservative with more strategies and longer blocks (\nz{v2b.E.size.spa_c_nb_ge.k100.b8}\% with a hundred strategies at block length 8, \nz{v2b.E.size.spa_c_nb_ge.k100.b22}\% at 22). Its size-adjusted power was close to the standard-deviation version's (\nz{v2b.F.adj.spa_c_nb_ge.iid_normal}\% against \nz{v2b.F.adj.spa_c_boot_t_ge.iid_normal}\% under i.i.d.\ normal returns), so what it loses is level, not ranking. A likely cause, which we have not tested, is that the observed and the resampled statistics are now studentized by different estimators: the natural block variance of a resample, computed from blocks of random length, is noisier than the sample's Politis--Romano estimate, which spreads the bootstrap distribution of the maximum and raises its critical value. A block studentization that treats the observed and the resampled statistics alike remains to be specified and tested; \citet{goncalves2011block} is a natural starting point.

\subsection{Where studentizing and recentring pay off}
\label{app:pay}
Studentizing exists for a reason, and experiment C2 shows when it helps and when it hurts (Table~\ref{tab:fair2}). With volatilities from 0.5 to 2, the strategy with the largest mean was the skilled one in \nz{v2b.C2.best_mean.low}\% of searches when the skilled strategy was the least volatile, \nz{v2b.C2.best_mean.median}\% when it had the median volatility and \nz{v2b.C2.best_mean.high}\% when it was the most volatile; the strategy with the largest $t$-statistic was the skilled one in about four searches in five wherever it sat (\nz{v2b.C2.best_t.low}\%, \nz{v2b.C2.best_t.median}\% and \nz{v2b.C2.best_t.high}\%). The Reality Check, which ranks by mean, accordingly rejected in \nz{v2b.C2.power.rc.low}\%, \nz{v2b.C2.power.rc.median}\% and \nz{v2b.C2.power.rc.high}\% of searches with a Sharpe ratio of 2, and the bootstrap-$t$ in \nz{v2b.C2.power.spa_c_boot_t_ge.low}\%, \nz{v2b.C2.power.spa_c_boot_t_ge.median}\% and \nz{v2b.C2.power.spa_c_boot_t_ge.high}\%: studentizing pays when the skilled strategy is among the less volatile and costs power when it is the most volatile (P15 \nz{v2b.P15.held}). The gain comes from ranking by $t$-statistic, not from the bootstrap: LET, the \v{S}id\'ak $t$ test, which does not resample, had about the bootstrap-$t$'s power (\nz{v2b.C2.power.sidak_t.low}\%, \nz{v2b.C2.power.sidak_t.median}\% and \nz{v2b.C2.power.sidak_t.high}\%). The fixed-variance version was again oversized (\nz{v2b.C2.size.spa_c_ge.het_null}\%), and the natural-block version again conservative (\nz{v2b.C2.size.spa_c_nb_ge.het_null}\%).

When the other strategies are clearly poor, the Reality Check, which recentres every strategy at its mean, is very conservative (size \nz{v2b.C2.size.rc.poor}\%, power \nz{v2b.C2.power.rc.poor}\%), and Hansen's consistent recentring raises power at the nominal level to \nz{v2b.C2.power.spa_c_ge.poor}\%, with or without studentizing, once ties are counted (sizes \nz{v2b.C2.size.spa_c_unstud.poor}\%, \nz{v2b.C2.size.spa_c_ge.poor}\% and \nz{v2b.C2.size.spa_c_boot_t_ge.poor}\%). Size-adjusted, the Reality Check ranks these searches as well as any test (\nz{v2b.C2.adj.rc.poor}\%): its loss is conservativeness at the nominal level, which is the point of \citet{hansen2005spa}. In the v2 version of this design (Table~\ref{tab:fair}), with the strict inequality, prediction P8 compared the raw power of an oversized bootstrap-$t$ (size \nz{v2.C.size.spa_c_boot_t.poor_alternatives}\%) with that of the conservative Reality Check; it held at the nominal level, not after size adjustment. Studentizing and recentring do what Hansen proposed them for. The lessons are to studentize inside the bootstrap, to count ties, and to remember that studentizing changes which strategies a test can find.
